\documentclass{article}

%-------------------------%
% PREAMBUŁA --------------%
%-------------------------%
\usepackage[utf8]{inputenc}
\usepackage[OT4]{polski}
\usepackage{caption}
\usepackage{amsmath}
\usepackage{tikz}
\usepackage{subcaption}
\usepackage{tabularx}
\usepackage{caption}
\usepackage{array}
\usepackage{hyperref}
\usepackage{graphicx}
\usepackage{fancyhdr}
\usepackage[justification=centering]{caption}
\usepackage[margin=60pt]{geometry}
\usepackage{longtable}



%\usepackage[T1]{fontenc}
%\usepackage{amsthm}
%\theoremstyle{plain}
%\usepackage{enumitem}
%\captionsetup[table]{skip=-4pt}
%\usepackage{tabulary}
%\usepackage{sidecap}
%\usepackage{gensymb}
%\usepackage{wrapfig}
%\usepackage{lastpage}
%\usepackage{gensymb}
%\usepackage{lipsum}
%\pagestyle{fancy}


\newlength{\RoundedBoxWidth}
\newsavebox{\GrayRoundedBox}
\newenvironment{GrayBox}[1][\dimexpr\textwidth-4.5ex]%
   {\setlength{\RoundedBoxWidth}{\dimexpr#1}
    \begin{lrbox}{\GrayRoundedBox}
       \begin{minipage}{\RoundedBoxWidth}}%
   {   \end{minipage}
    \end{lrbox}
    \begin{center}
    \begin{tikzpicture}%
       \draw node[draw=black,fill=black!1,rounded corners,%
             inner sep=2ex,text width=\RoundedBoxWidth]%
             {\usebox{\GrayRoundedBox}};
    \end{tikzpicture}
    \end{center}}

%------- Typy kolumn do ustawiania szerokości ------- %
\newcolumntype{L}[1]{>{\raggedright\let\newline\\\arraybackslash\hspace{0pt}}m{#1}}
\newcolumntype{C}[1]{>{\centering\let\newline\\\arraybackslash\hspace{0pt}}m{#1}}
\newcolumntype{R}[1]{>{\raggedleft\let\newline\\\arraybackslash\hspace{0pt}}m{#1}}








\begin{document}

%-------------------------%
%Tabela nagłówkowa -------%
%-------------------------%

\begin{table}[h]
\begin{tabular}{|l|l|l|l|l|l|}
\hline
\begin{tabular}[c]{@{}l@{}}
Wydział:\\ WFiIS\end{tabular} &
\multicolumn{2}{l|}{\begin{tabular}[c]{@{}l@{}}
Imię i nazwisko:\\ 1. Axel Zuziak\\ 2. Marcin Węglarz \end{tabular}} 
& Rok \textbf{II}         
& Grupa \textbf{B}          
& Zespół \textbf{03}      \\ \hline
\textbf{\begin{tabular}[c]{@{}l@{}}
LABOLATORIUM\\ TECHNIK \\ JĄDROWYCH\end{tabular}} &
\multicolumn{4}{l|}{ Temat:\textbf{ Rentgenowska analiza fluorescencyjna.}
	 }                                                                                                                       & \multicolumn{1}{c|}{\begin{tabular}[c]{@{}c@{}}
Nr ćwiczenia\\ \textbf{14}
\end{tabular}} \\ \hline
\begin{tabular}[c]{@{}c@{}}
Data wykonania:\\ 06.05.2015
\end{tabular} &
\begin{tabular}[c]{@{}c@{}}
 Data oddania:\\ 20.05.2015
\end{tabular} &
\begin{tabular}[c]{@{}c@{}}
Zwrot do poprawy: \\ 
\end{tabular}                                 
& 
\begin{tabular}[c]{@{}c@{}}
Data oddania: \\ 
\end{tabular}&
\begin{tabular}[c]{@{}c@{}}
Data zaliczenia:\\
\end{tabular}& 
\begin{tabular}[c]{@{}c@{}}
OCENA:\\
\end{tabular}       
\\ & & & & & \\ \hline
\end{tabular}
\end{table}
%-----Koniec tabeli------

%%%%%%%%%%%% Wstęp %%%%%%%%%%%%
\section{Wstęp}
Celem ćwiczenia było zapoznanie się z działaniem licznika proporcjonalnego oraz z podstawami analizy fluorescencyjnej.\\\\

\subsection{Licznik proporcjonalny}
Licznik proporcjonalny jest gazowym detektorem promieniowania jonizującego, w którym energia padającego promieniowania jest proporcjonalna do wyjściowego impulsu eletrycznego.\\
 Kluczową rolę w tego rodzaju detektorach pełni napięcie zasilania. Musi być ono tak dobrane, aby elektrony, powstałe wskutek jonizacji pierwotnej, przyspieszane polem $E$ miały wystarczającą energię do dalszej jonizacji gazu detektora. Zjawisko to nosi nazwę jonizacji wtórnej i prowadzi do powielenia jonów absorbowanych na elektrodach (wzmocnienie gazowe). Ponieważ jonizacja wtórna zachodzi blisko anody, w praktyce wartość współczynnika wzmocnienia gazowego nie zależy od miejsca jonizacji gazu detektora.\\\\
 Z powyższych rozważań możemy zapisać następującą proporcję:
 \begin{equation}
 U_{w} \propto  \frac{kM\Delta E}{W}
 \label{wz_U_w}
 \end{equation}
Gdzie:\\
$U_w$ - impuls po wzmocnieniu, trafiający do analizatora\\
$k$ - współczynnik wzmocnienia impulsu wychodzącego z detektora\\
$M$ - współczynnik wzmocnienia gazowego\\
$\Delta E$ - część energii promieniowania padającego deponowana w objętości czynnej detektora\\
$W$ - średnia wartość energii potrzebnej do wytworzenia jednej pary jonów.\\\\

\subsection{Zjawisko fluorescencji}
W przypadku padania promieniowania $\gamma$ na atom danego pierwiastka, możliwe jest wybicie (foto)elektronu z wewnętrznych powłok ($K$,$L$). W wyniku tego zjawiska elektron z wyższej powłoki dążąc do stanu o najniższej energii zapełni powstałą dziurę, emitując jednocześnie kwant promieniowania X o energii równej różnicy energii pomiędzy powłokami.\\
 Promieniowanie to nosi nazwę promieniowania charakterystycznego i jest podstawą analizy fluorescencyjnej. Nazwa tego promieniowania bierze się z faktu, że jest to promieniowanie charakterystyczne dla danego pierwiastka (każdy pierwiastek ma inne), dzięki czemu badając widmo takiego promieniowania jesteśmy w stanie określić jaki pierwiastek badamy.\\
 Zjawiskiem konkurencyjnym do promieniowania charakterystycznego jest emisja elektronów Augera. Jeżeli różnica pomiędzy progami energetycznymi powłok jest większa od energii progowej powłoki na którą elektron chce "wskoczyć", to może on zostać wyemitowany.\\
 W przeprowadzanym doświadczeniu źródłem kwantów gamma wywołujących zjawisko fluorescencji był $^{238}$Pu.
 
 \section{Przebieg doświadczenia}
 
\subsection{Wyznaczanie wzmocnienia gazowego}
W pierwszej części doświadczenia wyznaczano współczynnik wzmocnienia gazowego dla napięcia zasilania równego 1500V. \\
W tym celu do aparatury włożono źródło $^{238}$Pu promieniujące jedynie do góry i przykryto je płytką wykonaną z żelaza. Promieniowanie charakterystyczne powstające w żelaznej tarczy było następnie rejestrowane przez detektor umieszczony pod źródłem neutronowym (wymiary źródła były mniejsze od płytki, aby "nie zasłaniać" detektora). Następnie zmniejszano napięcie zasilania, aż do momentu, w którym pik (będący promieniowaniem charakterystycznym dla żelaza) pokrywał się z szumem.\\\\
Centrum piku dla wartości napięcia 1500V przypadało dla wartości 4V w wykorzystywanym analizatorze wielokanałowym. Przy każdorazowym zmniejszaniu napięcia detektora regulowano, przy pomocy współczynnika $k$ wzmocnienia wzmacniacza (sygnału wyjściowego z detektora), tak, aby centrum piku nie zmieniało swojego położenia (4V). Wyniki przedstawiono w tabeli \ref{tabela_wspl_wzmocnienia}. \\\\
W chwili, gdy szum pokrywa się z naszym pikiem przyjmujemy, że wzmocnienie gazowe detektora wynosi $M=1$. Ponieważ wartość $U_w$ pozostaje cały czas stała, korzystając ze wzoru (\ref{wz_U_w}) możemy napisać:\\
\[k_{1500}M_{1500} = k_{1200}M_{1200}
\]\\
W powyższej równości uwzględniono fakt, że szum pokrył się z badanym pikiem dla napięcia detektora równego 1200V. Stąd:\\
\[M_{1500} = \frac{k_{1200}M_{1200}}{k_{1500}} = \frac{1395.2 \cdot 1}{56.8} \approx 25
\]


\begin{table}[h!]
	\centering
	\caption{Tabela zawiera wartości napięcia detektora, a także odpowiadające im wartości wzmocnienia wzmacniacza $k$.}
	\begin{tabular}{|c|c|c|c|}
		\hline
		Lp. & Napięcie $U_d$ [V] & Współczynnik $k$ & Stosunek $\frac{k_i}{k_0}$ \\\hline
		1 & 1500 & 56.8 & 1 \\\hline
		2 & 1480 & 79.2 & 1.39 \\\hline
		3 & 1460 & 88.8 & 1.56 \\\hline
		4 & 1440 & 111.2 & 1.96 \\\hline
		5 & 1400 & 166.4 & 2.93 \\\hline
		6 & 1360 & 281.6 & 4.96 \\\hline
		7 & 1320 & 384 & 6.76 \\\hline
		8 & 1280 & 576 & 10.14 \\\hline
		9 & 1240 & 832 & 14.65 \\\hline
		10 & 1220 & 1190.4 & 20.96 \\\hline
		11 & 1200 & 1395.2 & 24.56 \\\hline
	\end{tabular}
\end{table}


Powyższa metoda nie jest metodą dokładną, ponieważ założenie, że dla napięcia 1200V $M$=1 jest przybliżeniem. Pozwala ona jednak na wyznaczenie rzędu wielkości wzmocnienia gazowego w detektorze proporcjonalnym.

\subsection{have fun :)}


%%%%%%%%%%%% BIBLIOGRAFIA %%%%%%%%%%%%
\newpage
\begin{thebibliography}{9}
	
	
	\bibitem{1}
	\url{http://nucleardata.nuclear.lu.se/toi/}

	
	\bibitem{DzK}
	B. Dziunikowski, S.J. Kalita
	\emph{Ćwiczenia laboratoryjne z jądrowych metod pomiarowcyh}, Wydawnictwa AGH, Kraków 1995
	
	%\bibitem{przybycien}
	%Mariusz Przybycień \emph{Tablice Statystyczne}, dostęp on-line\\
	%\url{http://home.agh.edu.pl/~mariuszp/wfiis_stat/tablice_ps_wir.pdf}
\end{thebibliography}
\vspace{2cm}








\end{document}





%\textsc{http://nucleardata.nuclear.lu.se/toi/}
%do bibliografii






